\section{Limitations and Responsible Use}
\label{sec:limitations}

\textbf{Reference design and label scope.} Gold150 is a disagreement-selected reference used in handbook development. The separate 50-sentence agreement study tests use of the final handbook; it does not supply a new held-out model test. Human annotation agreement and model scores use different samples and targets. K includes qualification and industry-background proxies, and L has only two reference spans. Independent annotation of new, advertisement-isolated texts and human review of K subtypes remain necessary to establish broader validity.

\textbf{Historical supervision and access.} The reference, original supervision, and agreement study use different handbook versions. Their labels remain frozen, so final-handbook consistency does not establish compatibility of all historical targets. The original corpus has disjoint source-document partitions, and later B2 checks exclude reference sentences from supervision. The later experimental splits nevertheless share source-document groups, which limits claims about new-document generalization. Incomplete pretraining records leave additional exposure uncertain. Public labels and predictions support inspection and rescoring, while complete expanded-pool inputs are not publicly available for retraining. Reuse must respect the component-specific source permissions.

\textbf{Experimental scope.} The Qwen comparison measures end-to-end gains from adaptation, including improved output validity. It does not isolate a semantic-extraction effect. Prompted configurations differ in inference budgets and access routes, expanded-pool studies use single seeds, and three-seed SD is imprecisely estimated. Source-group bootstrap intervals address dependence among recorded reference prefixes but remain conditional on the selected examples and trained models. Eight empty-reference sentences support descriptive error counts only.

\section{Conclusion}
\label{sec:conclusion}

Chinese-SkillSpan provides a benchmark for delimiting and typing competency mentions in Chinese job advertisements. Its contributions are an original-text resource with distinct human and Silver labels, an ESCO-informed operational handbook for Chinese boundaries and contextual types, and annotation-quality evidence that separates independent agreement from assisted review. Fixed character offsets, released predictions, and a shared scorer make model errors inspectable.

Three independent annotators achieved mean typed exact agreement of 0.818 after individual self-review under the final handbook. Qwen adaptation improved end-to-end extraction against its own base checkpoint, while the prompted GPT configuration remained strongest by point score. The Chinese encoder comparison did not establish an ESCO initialization advantage. Together, these experiments demonstrate use of the resource for diagnostic comparison and identify the additional annotation and independent testing needed for broader deployment.
